\section{Chellas's standard normal systems}
\label{sec:p08-chellas-standard}

The standard-model presentation in Chellas's 1980 book starts with a
relational semantics, defines normal systems by closure under a finite
boxed-consequence operation, and obtains determination and finite
countermodels by different constructions. The following treatment keeps
those three interfaces distinct. Its source is Chapters 3--5, with the
language, deducibility and maximality conventions of Chapter 2
\cite[chs.~2--5]{Chellas1980}. The systems treated here are
\(K,KD,KT,KB,KTB,K4,S4,S4.2,S5\). In particular, the symmetry-only
system \(KB\) is not the Brouwersche system \(KTB\).

\subsection{The source language, theorem systems and standard models}

\begin{hbconvention}{def:p08-chellas-source-language}{Chellas's primitive language}
The source's sentences have distinct atoms \(P_0,P_1,\ldots\), constants
\(\topv,\bottom\), Boolean connectives
\(\neg,\land,\lor,\to,\leftrightarrow\), and both modal operators
\(\Box,\Diamond\). They are finite, uniquely readable expressions.
Here \(\operatorname{Sn}(A)\) is their finite set of subsentences,
including \(A\) itself. Superscripts on modalities count iterations:
\(\Box^0 A=\Diamond^0 A=A\).

Unlike the primitive-box sheets, this syntax has primitive diamond
sentences. The schema \(\mathrm{Df}\Diamond\) is
\(\Diamond A\leftrightarrow\neg\Box\neg A\); it is a theorem
schema, not a claim that these two source expressions are syntactically
identical. The source's truth and falsity glyphs are rendered by the
handbook's \(\topv,\bottom\). Its uppercase atom indices and both
primitive modalities are retained.
\end{hbconvention}

\begin{hbdefinition}{def:p08-chellas-source-consequence}{Normal systems and source deducibility}
A source system \(\Sigma\) is a set of sentences closed under
tautological consequence. It is \emph{normal} when it contains every
\(\mathrm{Df}\Diamond\) instance and, for every finite list
\(A_1,\ldots,A_r\), its theoremhood is closed under the operation
\[
 (A_1\land\cdots\land A_r)\to B\in\Sigma
 \quad\Longrightarrow\quad
 (\Box A_1\land\cdots\land\Box A_r)\to\Box B\in\Sigma.
\]
For the empty list, these expressions mean \(B\) and \(\Box B\).
Write \(\vdash_\Sigma A\) for \(A\in\Sigma\), and define
\[
 \Gamma\vdash_\Sigma A
 \quad\Longleftrightarrow\quad
 \text{some }A_1,\ldots,A_r\in\Gamma\text{ have }
 (A_1\land\cdots\land A_r)\to A\in\Sigma.
\]
The premise set may be arbitrary. A set is \(\Sigma\)-consistent when
it does not deduce \(\bottom\); maximality is maximal consistent
extension in the full countable source language. These are Chellas's
finitely supported deducibility and theoremhood conventions, not a
global-premise rule that permits boxing an assumption
\cite[defs.~2.11--2.17, 4.1]{Chellas1980}.
\end{hbdefinition}

\begin{hbdefinition}{def:p08-chellas-standard-model}{Source standard models and degrees of validity}
A standard model is \(\mathcal M=\langle W,R,P\rangle\), with
\(R\subseteq W^2\) and \(P:\N\to\mathcal P(W)\).
The source allows empty \(W\); a pointed model necessarily has a
world and hence nonempty \(W\). The typed use of \(P_n\) in
\(\alpha\in P_n\) denotes the valuation set \(P(n)\), whereas
\(P_n\) in a sentence denotes its atom. Truth is written
\(\models_\alpha^{\mathcal M}A\), and has the classical Boolean
clauses and
\[
\begin{aligned}
 \models_\alpha^{\mathcal M}P_n&\ \Longleftrightarrow\ \alpha\in P_n,\\
 \models_\alpha^{\mathcal M}\Box A&\ \Longleftrightarrow\
       \forall\beta\,(\alpha R\beta\Rightarrow
                     \models_\beta^{\mathcal M}A),\\
 \models_\alpha^{\mathcal M}\Diamond A&\ \Longleftrightarrow\
       \exists\beta\,(\alpha R\beta\land
                     \models_\beta^{\mathcal M}A).
\end{aligned}
\]
The truth set is \(\|A\|^{\mathcal M}=\{\alpha\in W:
\models_\alpha^{\mathcal M}A\}\). Truth in a model means truth at
every world; validity in a model class \(C\), written \(\models_C A\),
means truth in every member. A system is \emph{determined} by \(C\)
when its theorems are exactly the sentences valid in \(C\).
Local consequence \(\Gamma\models_C^\ell A\) instead tests a
single world satisfying every premise. Empty models do not alter either
theoremhood validity or local consequence on a class that also contains
the stated nonempty models.
\end{hbdefinition}

\begin{hblemma}{thm:p08-chellas-normal-interface}{proved=annex:p8.chellas.normal-interface}{Normal closure and the registered Hilbert presentation}
Normal closure is equivalent to propositional logic, distribution K,
\(\mathrm{Df}\Diamond\), modus ponens and theorem-only
necessitation. It permits replacement of theorem-equivalent expressions
inside any sentence, including under either modality. Eliminating
\(\Diamond,\leftrightarrow,\topv\) and replacing source atoms
\(P_n\) by the sheet atoms \(p_n\) preserves theoremhood,
source deducibility and truth. Thus each selected source system has
the theorem-only local-premise interface fixed here and, for the eight
systems with the sheets listed below, exactly that sheet's theory. The \(KB\)
extension consists of the B schema alone over K; \(KTB\) also includes
T. The \(S4.2\) extension is \(KT4G\), where G is .2.
\end{hblemma}

The registered surfaces for K, KD, KT, KTB, K4, S4, S4.2 and S5
are Sheets~\ref{sheet:hilbert-k}, \ref{sheet:hilbert-d},
\ref{sheet:hilbert-t}, \ref{sheet:hilbert-b}, \ref{sheet:hilbert-k4},
\ref{sheet:hilbert-s4}, \ref{sheet:hilbert-s42} and
\ref{sheet:hilbert-s5}. The explicit KB specification above has no T
component; it is not silently assigned the KTB sheet.

\subsection{Axioms, frames and the selected inclusion order}

\begin{hbdefinition}{def:p08-chellas-characteristic-schemas}{Characteristic schemas and their source names}
The source names D, T, B, 4 and 5 designate, respectively,
\[
 \Box A\to\Diamond A,\quad \Box A\to A,\quad
 A\to\Box\Diamond A,\quad \Box A\to\Box\Box A,\quad
 \Diamond A\to\Box\Diamond A.
\]
G designates \(\Diamond\Box A\to\Box\Diamond A\).
Its frame condition is directed convergence above a world, as in
Definition~\ref{def:p08-modal-frame-properties}. The source calls that
condition ``incestuality.'' It differs from \emph{global} convergence:
\(\forall\beta\forall\gamma\exists\delta\,
 (\beta R\delta\land\gamma R\delta)\).
For a reflexive transitive frame generated from a root, the two conditions
coincide. This distinction is needed in the S4.2 finite-model construction.
\end{hbdefinition}

\begin{table}[!htbp]
\centering\small
\begin{tabular}{@{}lll@{}}
\toprule
Source system & Added schemas over K & Standard frame class \\
\midrule
K & none & arbitrary \\
KD & D & serial \\
KT & T & reflexive \\
KB & B & symmetric \\
KTB & T, B & reflexive, symmetric \\
K4 & 4 & transitive \\
S4 = KT4 & T, 4 & reflexive, transitive \\
S4.2 = KT4G & T, 4, G & reflexive, transitive, locally directed \\
S5 = KT5 & T, 5 & equivalence \\
\bottomrule
\end{tabular}
\caption{Chellas's selected standard systems. These are exact axiom/frame
pairings, not an assertion that an arbitrary normal extension is complete
for the frame conditions of whichever axioms it happens to contain.}
\label{tab:p08-chellas-standard-systems}
\end{table}

\begin{hblemma}{thm:p08-chellas-correspondence}{proved=annex:p8.chellas.correspondence}{All-valuations correspondences and soundness}
On a fixed frame, validity of every instance of D, T, B, 4, 5 or G
is equivalent, respectively, to seriality, reflexivity, symmetry,
transitivity, Euclideanness or local directed convergence. Each row of
Table~\ref{tab:p08-chellas-standard-systems} is sound for local
consequence under source deducibility. Reflexivity and Euclideanness
together make an equivalence relation.
\end{hblemma}

\begin{hbtheorem}{thm:p08-chellas-selected-inclusions}{proved=annex:p8.chellas.selected-inclusions}{Exact inclusion order on the nine displayed systems}
The generating strict inclusions among these nine systems are
\[
\begin{gathered}
 K\subsetneq KD\subsetneq KT\subsetneq KTB\subsetneq S5,\qquad
 K\subsetneq KB\subsetneq KTB,\\
 K\subsetneq K4\subsetneq S4\subsetneq S4.2\subsetneq S5,
 \qquad KT\subsetneq S4.
\end{gathered}
\]
Their transitive consequences are all the inclusions in this finite
selection; every other pair is incomparable. This is an induced order on
nine named logics, not the lattice of all normal modal logics. In
particular, the B schema, KB and KTB are three different items of information.
\end{hbtheorem}

The source also obtains alternative axiomatizations. T entails D, and
over a normal B-system, 4 and 5 entail each other. Thus
\(KT5=KTB4=KDB4=KDB5\). The point is the independent contribution of
each schema: B supplies symmetry, 4 supplies transitivity, and T supplies
reflexivity. Annex~\ref{annex:p8.chellas.selected-inclusions} gives the
syntactic derivations and complete finite separating models
\cite[\S4.3, thm.~5.2]{Chellas1980}.

\subsection{Canonical determination and the G family}

\begin{hbdefinition}{def:p08-chellas-proper-canonical}{Proper canonical standard model}
For a consistent normal system \(\Sigma\), let \(W_\Sigma\) be all
maximal \(\Sigma\)-consistent sets and put
\[
 \alpha R_\Sigma\beta\ \Longleftrightarrow\
       \{A:\Box A\in\alpha\}\subseteq\beta,
 \qquad P_n=\{\alpha\in W_\Sigma:P_n\in\alpha\}.
\]
The source calls this the \emph{proper} canonical standard model.
Its more general definition of a canonical standard model only requires
\(\Box A\in\alpha\) exactly when all its successors contain \(A\);
the displayed relation is a specific construction satisfying that
requirement, not merely that requirement repeated as a definition
\cite[defs.~5.5, 5.9]{Chellas1980}.
\end{hbdefinition}

\begin{hblemma}{thm:p08-chellas-source-canonical}{proved=annex:p8.chellas.source-canonical}{The source canonical construction for a normal system}
Every consistent set in a consistent normal system on the source's
countable language has a maximal extension. The proper canonical model
is nonempty and satisfies
\[
 \models_\alpha^{\mathcal M_\Sigma}A\ \Longleftrightarrow\ A\in\alpha,
 \qquad
 \alpha R_\Sigma\beta\ \Longleftrightarrow\
       \{\Diamond A:A\in\beta\}\subseteq\alpha.
\]
It follows that \(\Sigma\) is determined by this particular model.
Completeness for a specified frame class additionally requires proving
that this model's frame belongs to that class.
\end{hblemma}

The seven previously completed canonical applications are reusable
without changing their consequence relation:
Theorem~\ref{thm:p08-modal-canonical-completeness} covers K, KT, S4
and S5, and Theorems~\ref{thm:p08-kd-strong-completeness},
\ref{thm:p08-ktb-strong-completeness} and
\ref{thm:p08-k4-strong-completeness} cover KD, KTB and K4.
Lemma~\ref{thm:p08-chellas-normal-interface} supplies the explicit
source-to-sheet bridge. The source canonical lemmas above and the two
following applications also give the missing source interfaces for KB
and S4.2; none relies on an assumed finite-model theorem.

\begin{hbtheorem}{thm:p08-chellas-kb-completeness}{proved=annex:p8.chellas.kb-completeness}{KB: strong local completeness without T}
For every source premise set \(\Gamma\) and source sentence \(A\),
\(\Gamma\vdash_{KB}A\) exactly when
\(\Gamma\models_{\mathrm{Sym}}^\ell A\), where
\(\mathrm{Sym}\) is the class of symmetric standard frames.
The KB canonical relation is symmetric, but need not be reflexive,
serial or transitive. KB is strictly weaker than KTB.
\end{hbtheorem}

\begin{hbdefinition}{def:p08-chellas-g-family}{Source G indices and path relations}
Define \(R^0\) as identity and \(\alpha R^{j+1}\beta\) when some
\(\gamma\) satisfies \(\alpha R\gamma R^j\beta\).
Chellas's four-index schema is
\[
 G^{k,l,m,n}:\quad
 \Diamond^k\Box^l A\to\Box^m\Diamond^n A,
 \qquad k,l,m,n\in\N.
\]
Its frame condition is
\[
 \alpha R^k\beta\land\alpha R^m\gamma
 \quad\Longrightarrow\quad
 \exists\delta\,(\beta R^l\delta\land\gamma R^n\delta).
\]
Zero indices are genuine identity cases, not positive-path abbreviations.
The characteristic schemas are
\(D=G^{0,1,0,1}\), \(T=G^{0,1,0,0}\),
\(B=G^{0,0,1,1}\), \(4=G^{0,1,2,0}\),
\(5=G^{1,0,1,1}\), and \(G=G^{1,1,1,1}\)
\cite[\S3.3]{Chellas1980}.
\end{hbdefinition}

\begin{hbtheorem}{thm:p08-chellas-g-determination}{proved=annex:p8.chellas.g-determination}{Full source G-family determination}
For every four nonnegative indices, the least normal system containing
\(G^{k,l,m,n}\) is determined by the displayed four-index frame
condition and is strongly complete for its local consequence.
In any normal extension containing that schema, the proper canonical
frame has the same condition. The latter fact does not imply that an
arbitrary stronger extension is sound for every frame satisfying just
this one condition.
\end{hbtheorem}

\begin{hbtheorem}{thm:p08-chellas-s42-completeness}{proved=annex:p8.chellas.s42-completeness}{S4.2: local and generated-convergent determination}
S4.2 is strongly complete for local consequence on reflexive transitive
locally directed frames. It is also strongly complete on reflexive
transitive \emph{globally} convergent frames. A pointed countermodel
for the first claim restricts to the submodel generated by its evaluation
world, and that restriction has the second frame property while
preserving every sentence and every premise at that world.
\end{hbtheorem}

\subsection{Generated models and both filtration constructions}

\begin{hbdefinition}{def:p08-chellas-generated-model}{Generated restriction}
For a world \(\alpha\) of \(\mathcal M\), let
\(W^\alpha=\{\beta:\alpha R^j\beta\text{ for some }j\geq0\}\).
Restrict \(R\) and every \(P_n\) to this set to obtain
\(\mathcal M^\alpha\). No outgoing edge from \(W^\alpha\) is
discarded. In an equivalence frame this generated frame is universal;
in a reflexive transitive frame it consists of all successors of its root.
\end{hbdefinition}

\begin{hblemma}{thm:p08-chellas-generated-truth}{proved=annex:p8.chellas.generated-truth}{Truth and validity under generated restrictions}
Every sentence has the same truth value at every retained world in
\(\mathcal M\) and \(\mathcal M^\alpha\). Thus a class of models
and the class of all its generated restrictions have the same validities.
In particular, S5 is determined by universal models as well as equivalence
frames; their finite counterparts determine the same theoremhood too.
\end{hblemma}

\begin{hbdefinition}{def:p08-chellas-source-filtration}{The source's two-modal filtration conditions}
Let \(\Gamma\) be closed under subsentences. In this subsection it
indexes filtration sentences, not premises. Worlds agree under
\(\alpha\sim_\Gamma\beta\) when they agree on every member of
\(\Gamma\); let \([\alpha]\) be the resulting class and
\([X]=\{[\alpha]:\alpha\in X\}\).
A source filtration \(\mathcal M^*=\langle W^*,R^*,P^*\rangle\)
has \(W^*=[W]\), \(P_n^*=[P_n]\) for atoms in \(\Gamma\), and
the following three relational requirements:
\[
\begin{aligned}
 \alpha R\beta&\Rightarrow[\alpha]R^*[\beta],\\
 [\alpha]R^*[\beta],\ \Box A\in\Gamma,\
       \models_\alpha^{\mathcal M}\Box A
       &\Rightarrow\models_\beta^{\mathcal M}A,\\
 [\alpha]R^*[\beta],\ \Diamond A\in\Gamma,\
       \models_\beta^{\mathcal M}A
       &\Rightarrow\models_\alpha^{\mathcal M}\Diamond A.
\end{aligned}
\]
The \emph{finest} relation relates classes when some representatives are
R-related; the \emph{coarsest} relation contains exactly the pairs
meeting the last two requirements. Any relation between these two is
a filtration relation. Both modal constraints are necessary because
the source language has both primitive modalities
\cite[def.~3.14, \S3.5]{Chellas1980}.
\end{hbdefinition}

\begin{hbtheorem}{thm:p08-chellas-filtration-truth}{proved=annex:p8.chellas.filtration-truth}{Full filtration truth, Boolean extension and size}
For every source filtration and \(A\in\Gamma\),
\(\models_\alpha^{\mathcal M}A\) exactly when
\(\models_{[\alpha]}^{\mathcal M^*}A\).
This also holds for the Boolean closure of \(\Gamma\).
If \(\Gamma\) has \(s\) sentences, the filtration has at most
\(2^s\) worlds. More generally, if every sentence of \(\Gamma\)
is model-equivalent to one of \(s\) sentences, the same bound holds.
Every filtration of a serial, reflexive or universal model retains that
property. Symmetry and transitivity require the further constructions below.
\end{hbtheorem}

\begin{hbdefinition}{def:p08-chellas-filtration-conditions}{Chellas's paired conditions \(c_1,c_2,c_3,c_4\)}
Let \(c_1(\alpha,\beta)\) denote the two modal preservation
implications in Definition~\ref{def:p08-chellas-source-filtration}.
Put \(c_2(\alpha,\beta)=c_1(\beta,\alpha)\).
The remaining source conditions are
\[
\begin{aligned}
 c_3(\alpha,\beta):\quad &
  \Box A\in\Gamma,\ \models_\alpha^{\mathcal M}\Box A
       \Rightarrow\models_\beta^{\mathcal M}\Box A,\\[-2pt]
 &\Diamond A\in\Gamma,\ \models_\beta^{\mathcal M}\Diamond A
       \Rightarrow\models_\alpha^{\mathcal M}\Diamond A;\\
 c_4(\alpha,\beta):\quad &c_3(\beta,\alpha).
\end{aligned}
\]
Every displayed universal constraint quantifies over the appropriate
modal sentences in \(\Gamma\). Because \(\Gamma\) contains their
matrices, these conditions depend on classes, not chosen representatives.
\end{hbdefinition}

\begin{hbtheorem}{thm:p08-chellas-paired-filtrations}{proved=annex:p8.chellas.paired-filtrations}{All four paired-condition constructions}
Defining \(R^*\) by \(c_1\land c_2\) produces a symmetric
filtration of any symmetric model. Defining it by \(c_1\land c_3\)
produces a transitive filtration of any transitive model.
The conjunction \(c_1\land c_2\land c_3\land c_4\) produces a symmetric
transitive filtration of any symmetric transitive model.
Finally, \(c_1\land c_3\land c_4\) produces a transitive Euclidean
filtration of any transitive Euclidean model. Seriality and reflexivity,
where present in the original model, persist in all these filtrations.
\end{hbtheorem}

\begin{hbtheorem}{thm:p08-chellas-modal-closure-filtration}{proved=annex:p8.chellas.modal-closure-filtration}{The second route: coarsest filtration through modal closure}
If \(\Gamma\) is closed under all modalities and subsentences, its
coarsest filtration preserves symmetry, transitivity and Euclideanness
whenever the original model has the corresponding property.
Modal closure of a finite set can be infinite as syntax yet logically
finite in a model. In particular, a euclidean model admits at most
fourteen inequivalent modalities, so modal closure of \(s\) sentences
has at most \(14s\) truth-set representatives and its filtration is
finite. This auxiliary K5 construction explains the source's second
route; it does not make arbitrary finite filtrations Euclidean.
\end{hbtheorem}

\begin{hbtheorem}{thm:p08-chellas-filtration-boundaries}{proved=annex:p8.chellas.filtration-boundaries}{Why the added conditions cannot be dropped}
A finest filtration of a symmetric model is symmetric, but an arbitrary
filtration need not be. A finest filtration of a transitive model need
not be transitive. A filtration of a locally directed model need not
be locally directed. Nor does an arbitrary filtration of a Euclidean
model have to be Euclidean. Complete finite examples in the reciprocal
proof show each failure; the last also reproduces the source's
one-atom, two-sentence, five-world construction.
\end{hbtheorem}

\subsection{Finite countermodels, effective decisions and characteristic modalities}

\begin{hbtheorem}{thm:p08-chellas-finite-models}{proved=annex:p8.chellas.finite-models}{Finite determination of all nine selected systems}
Each selected system is determined by the finite models in its stated
frame class. Every non-theorem \(A\) has a pointed countermodel
with at most \(2^{|\operatorname{Sn}(A)|}\) worlds.
For S4.2, first generate from the countermodel world, then apply
the transitive paired-condition filtration: global convergence of the
generated model supplies convergence in the quotient. Adding edges to
an arbitrary filtered fork without checking truth is not this argument.
The finite-model statement concerns theoremhood and finite inputs, not
arbitrary infinite premise sets on finite frames.
\end{hbtheorem}

\begin{hbmethod}{def:p08-chellas-finite-search}{A complete finite countermodel search}
Given a source sentence \(A\), compute its atoms and subsentences and
set \(b=2^{|\operatorname{Sn}(A)|}\). For each \(1\leq j\leq b\),
enumerate all binary relations on \(\{0,\ldots,j-1\}\), retain the
ones satisfying the chosen row's first-order frame test, and enumerate
all subsets of that set for each atom occurring in \(A\).
Evaluate \(A\) bottom-up at every world. Return ``non-theorem''
with its frame, valuation and falsifying world when a failure appears;
if the finite enumeration ends without one, return ``theorem.''
This algorithm can consume enormous time and space; decidability is
not a claim of practical efficiency or a sharper complexity bound.
\end{hbmethod}

\begin{hbtheorem}{thm:p08-chellas-decision}{proved=annex:p8.chellas.decision}{Termination and exact finite-input decisions}
Method~\ref{def:p08-chellas-finite-search} terminates and decides
theoremhood for each selected system. For a finite list \(\Gamma_0\),
it decides local consequence by testing
\((\bigwedge\Gamma_0)\to A\), and satisfiability by testing whether
\(\neg A\) is a theorem. An alternative source procedure dovetails
effective proof enumeration with finite countermodel enumeration; one
of its two searches must terminate. No decidability claim is made for
an arbitrary infinitely specified premise set.
\end{hbtheorem}

\begin{hbtheorem}{thm:p08-chellas-modalities}{proved=annex:p8.chellas.modalities}{Exact characteristic modality counts in the selected systems}
K, KD, KT, KB, KTB and K4 each have infinitely many inequivalent
modalities. S4 has fourteen: the seven affirmative representatives
\[
 \text{identity},\ \Box,\ \Diamond,\ \Box\Diamond,
 \Diamond\Box,\ \Box\Diamond\Box,\ \Diamond\Box\Diamond
\]
and their negations. S4.2 has ten: the first five affirmative
representatives and their negations. S5 has six: identity, box,
diamond and their negations. Reduction identities establish upper
bounds; explicit finite truth-pattern witnesses establish distinctness.
These are counts for the named systems themselves, not for every
stronger normal extension \cite[\S4.4, \S5.2, exs.~4.66, 5.16]{Chellas1980}.
\end{hbtheorem}

\paragraph{Use and limits.}
Select the logic by the intended accessibility assumptions, not by the
informal word ``necessary.'' For example, seriality excludes a world with
no admissible alternatives, reflexivity makes necessities true at their
current world, and convergence requires compatible refinements of
alternatives. The same formal semantics can receive different intended
interpretations; frame completeness proves the displayed mathematical
consequence claim, not the adequacy of a philosophical interpretation.
GL, Grz, arithmetic realizations and neighborhood semantics are distinct
presentations and are not supplied by this standard-system construction.
